\chapter{Automated Fault Detection, Categorization, and Diagnosis for Transformer Architectures}
\label{chapter_defaultpp}

Chapter~\ref{chapter_default} showed that hierarchical diagnosis can detect and categorize faults in FFNNs, CNNs, and RNNs, but that result does not transfer directly to transformer architectures. Transformer faults are often silent. A stale QKV projection, a missing attention-score scaling factor, or a corrupted residual path can leave the model runnable, show a normal loss curve, and still degrade behavior without raising an exception or producing NaNs. The debugging problem is therefore not only to detect that something is wrong, but also to determine which transformer component is responsible and what mechanism caused the failure. This chapter studies that problem on the labeled benchmark produced by the FrankenFormer mutation-testing framework~\cite{frankenformer}. The dataset contains 3{,}739 averaged configurations across seven transformer models and nine tasks. On this benchmark, we address three linked tasks: Level~1 fault detection, Level~2 fault categorization into architecture-specific components, and Level~3 root-cause diagnosis within the predicted fault category.

We introduce DEFault++ (Detect, Explain, and Analyze Faults in Transformers), which extends DEFault's hierarchical reasoning strategy to attention-based architectures. Unlike prior DNN-level techniques, DEFault++ uses transformer-specific runtime measurements and a component-level label space derived from the attention fault taxonomy~\cite{attention_taxonomy}. It also organizes the diagnosis around a Fault Propagation Graph (FPG), a deterministic directed graph that encodes how faults propagate between transformer components. This design keeps the resulting prediction tied to transformer subsystems rather than to a flat feature vector.

The chapter proceeds from motivation and background (Sections~\ref{sec:defaultpp_introduction}--\ref{sec:defaultpp_background}) through the DEFault++ technique (Section~\ref{sec:defaultpp_approach}), experimental evaluation (Section~\ref{sec:defaultpp_experiment}), case study (Section~\ref{sec:defaultpp_casestudy}), threats to validity (Section~\ref{sec:defaultpp_threats}), related work (Section~\ref{sec:defaultpp_related}), and contributions (Section~\ref{sec:defaultpp_summary}).

\section{Motivation}
\label{sec:defaultpp_introduction}

Observable symptoms of transformer faults, such as degraded output quality, ineffective fine-tuning, or unexpected calibration drift, are rarely specific to one component. The same externally visible symptom may arise from projection errors, masking faults, positional faults, residual corruption, or broader training instabilities. A developer can inspect loss curves, gradients, attention maps, and parameter updates manually, but those metrics do not by themselves provide a structured path from symptom to architectural cause.

Existing DNN debugging techniques do not resolve this ambiguity. Rule-based approaches such as UMLAUT~\cite{umlaut} and AutoTrainer~\cite{autotrainer} encode common training symptoms, including exploding gradients, vanishing gradients, and oscillating loss. Learning-based approaches such as DeepFD~\cite{deepfd} improve over fixed rules by training on runtime features. Our prior work DEFault~\cite{default} further showed that hierarchical classification can improve detection and diagnosis across FFNNs, CNNs, and RNNs. However, all of these techniques, including DEFault, operate with either generic DNN symptom rules or architecture-agnostic features. They were not designed to represent projection alignment, attention-score behavior, masking integrity, or KV-cache dynamics at the granularity needed for transformer diagnosis.

Three results make a transformer-specific diagnostic technique feasible. First, Chapter~\ref{chapter_attention} established an empirical taxonomy of 555 attention-related faults and organized them into component-level categories with finer-grained root causes~\cite{attention_taxonomy}. Second, Hessian-based analysis showed that perturbations in attention mechanisms propagate across projections, heads, and residual paths rather than remaining confined to one local variable~\cite{hessian_analysis}. Third, the FrankenFormer mutation-testing framework produced a labeled benchmark with instrumented traces over multiple transformer models and tasks~\cite{frankenformer}. Together, these results supply a target label space, a structural view of propagation, and the training data needed for supervised diagnosis.

DEFault++ combines those three results. It preserves DEFault's hierarchical reasoning strategy but replaces the DNN-level feature space and label space with transformer-specific ones. It then adds a Level~3 diagnosis stage and component-aligned explanation through the FPG. The aim is not merely to predict that a transformer is faulty, but to narrow the fault to a fault category and then to a root cause that a developer can act on.

We evaluate DEFault++ on 3{,}739 instances across encoder and decoder architectures and compare it with four prior techniques on the directly comparable fault-detection task. The chapter makes the following contributions:

\begin{enumerate}[label=(\alph*)]
\item A novel hierarchical neural technique for fault detection, fault categorization, and root-cause diagnosis in transformer architectures.

\item A Fault Propagation Graph (FPG) derived from eight deterministic propagation mechanisms based on the forward and backward pass of transformer architectures. 

\item A custom root-cause separation loss for Level~3. It combines a contrastive term with a prototype-matching loss and improves root-cause diagnosis over the base variant.
\end{enumerate}

%==============================================================================
\section{Motivating Example}
\label{sec:defaultpp_motivation}

Listing~\ref{lst:qkv_fusion_bug} shows a real-world fault from the HuggingFace Diffusers library (Issue~\#11903\footnote{\href{https://github.com/huggingface/diffusers/issues/11903}{HuggingFace Diffusers Issue \#11903}}). The fault resides in a specific attention subsystem. The observable symptom at the task level gives no indication of which component is affected.

\begin{listing}[!htbp]
\tcbinputlisting{
    thesis-code-scriptsize,
    unbreakable,
    minted language=python,
    listing file={code/defaultplusplus/qkv_fusion_bug.py},
}
\caption{Faulty QKV projection fusion from HuggingFace Diffusers Issue \#11903 (simplified).}
\label{lst:qkv_fusion_bug}
\end{listing}

The \texttt{fuse\_qkv\_projections} method (line~12) combines separate \texttt{q\_proj}, \texttt{k\_proj}, and \texttt{v\_proj} layers into a single \texttt{to\_qkv} layer. After fusion, the \texttt{forward} method (line~21) routes all computation through \texttt{to\_qkv}, while \texttt{q\_proj} and \texttt{v\_proj} remain in memory with their original parameters. LoRA and EMA updates apply to \texttt{q\_proj} and \texttt{v\_proj}, bypassing \texttt{to\_qkv} entirely.

Training proceeds without exceptions. No NaN or Inf values appear. The loss curve decreases without anomaly. The only observable symptom is that fine-tuning produces no effect on model behavior. This symptom does not identify the affected component.

Existing techniques do not detect this fault. AutoTrainer~\cite{autotrainer} monitors generic training symptoms such as oscillating loss and vanishing gradients. These symptoms are absent here. UMLAUT~\cite{umlaut} applies predefined heuristics over validation trends. A plausible loss curve does not trigger any of its rules. DeepFD~\cite{deepfd} and DeepLocalize~\cite{deeplocalize} classify DNN-level fault categories and numerical anomalies. Neither examines projection-layer structure or update-path coherence. ATTNChecker~\cite{liang2025attnchecker} targets transient NaN/Inf-style failures in attention computations. The stale projection configuration produces neither.

DEFault++ instruments training to measure QKV alignment, attention entropy, and projection update activity. These measurements characterize the affected subsystem, not task-level behavior alone. DEFault++ detects the fault and assigns it to a component-level category. It then identifies the root cause within that category.

%==============================================================================
\section{Background}
\label{sec:defaultpp_background}

\textit{Fault taxonomies.} Prior empirical studies identified training and model faults as the dominant fault types in DNN programs and organized them into taxonomies by root cause~\cite{faulttaxonomy,dlbugcharacterstics}. Those taxonomies cover weights, regularization, hyperparameters, and layer configuration across FFNNs, CNNs, and RNNs. They do not include transformer-specific fault patterns, because they predate the widespread adoption of attention-based architectures. We established an empirical taxonomy of 555 attention-related faults~\cite{attention_taxonomy} that groups faults by the transformer component they affect: QKV projections, masking logic, positional encoding, residual paths, and KV caching. Within each component group, root-cause mechanisms are distinguished. A projection-dimension mismatch, a stale projection update path, and direct weight corruption all belong to the QKV category yet produce different alignment, entropy, and gradient patterns. The taxonomy supplies the Level~2 fault categories and the Level~3 root-cause targets that DEFault++ predicts.

\textit{Mutation-based fault injection.} Supervised diagnosis requires labeled faulty instances. DeepMutation~\cite{deepmutation} introduced source-level and model-level mutation operators for feed-forward networks, covering a narrow fault type set. DeepCrime~\cite{deepcrime} extended mutation to six fault categories across model construction and training and introduced the killability criterion, which requires that injected faults produce measurable behavioral change. DEFault~\cite{default} added ten layer-level operators and RNN support, producing a benchmark of 14{,}652 mutant programs across FFNNs, CNNs, and RNNs. None of these frameworks targets transformer architectures. FrankenFormer~\cite{frankenformer} applies the attention fault taxonomy~\cite{attention_taxonomy} to generate and validate transformer fault configurations. It produces 3{,}739 averaged configurations with component-level fault labels across encoder and decoder architectures, which serve as the supervised training data for DEFault++.

\textit{DNN debugging techniques.} Rule-based debuggers AutoTrainer~\cite{autotrainer} and DeepDiagnosis~\cite{deepdiagnosis} monitor generic training symptoms including exploding gradients, vanishing gradients, and oscillating loss. DeepFD~\cite{deepfd} trains classifiers on runtime features to identify five common fault types, improving over fixed rules. DEFault~\cite{default} showed that hierarchical classification further raises detection Macro-F1 above 0.85 across 14{,}652 DNN programs. More recently, tools have targeted transformer architectures directly: ATTNChecker~\cite{liang2025attnchecker} detects transient NaN/Inf-style faults in attention computations, FT-Transformer~\cite{fttransformer} addresses hardware-level fault tolerance for vision transformers under soft errors, and AtPatch~\cite{atpatch2026} localizes attention-map anomalies at the patch level. These techniques, whether operating on generic DNN features or on one transformer fault surface, share a common limitation: none observes projection alignment, attention-score behavior, masking integrity, or KV-cache dynamics across a component-level label space, and none distinguishes root causes within a fault category. Applied to the FrankenFormer benchmark~\cite{frankenformer}, AutoTrainer and DeepDiagnosis acheive detection Macro-F1 of only 0.284 and 0.338 on encoders, confirming that generic training symptoms do not transfer to transformer fault diagnosis.

DEFault++ measures transformer-specific runtime behavior and predicts component-level fault categories and root causes through a three-level hierarchy, targeting the diagnostic granularity that existing tools do not cover.

%==============================================================================
\section{DEFault++}
\label{sec:defaultpp_approach}

DEFault++ takes as input the labeled set of 3{,}739 averaged configurations, where each instance records the instrumented behavioral metrics for one configuration averaged across five seed runs. The three-level hierarchy introduced in Section~\ref{sec:defaultpp_introduction} is realized as follows. Level~1 uses the full feature representation for binary fault detection. Only instances predicted as faulty proceed to Level~2, which assigns one of the architecture-specific fault categories (11 for encoders, 12 for decoders). Level~3 predicts the root cause within the Level~2 predicted category through prototypical classification in an embedding space aligned with the transformer's Fault Propagation Graph. The same embedding space also produces feature-group salience for each diagnosis. Figure~\ref{fig:defaultpp_pipeline} shows the three-level hierarchy.

\begin{figure}[!htbp]
\centering
\includegraphics[width=\linewidth]{fig_method_overview.pdf}
\caption{Overview of the DEFault++ technique.}
\label{fig:defaultpp_pipeline}
\end{figure}

%------------------------------------------------------------------------------
\subsection{Dataset Construction}
\label{sec:defaultpp_dataset}

We construct the input dataset using the mutation-testing framework from Chapter~\ref{chapter_frankenformer}~\cite{frankenformer}. Each fault configuration $\mathcal{C} = (m, t, u, f, v, \ell, \sigma)$ specifies a transformer model~$m$, a fine-tuning task~$t$, a target component~$u$, a fault category~$f$, a fault variant~$v$, a target layer~$\ell$, and a severity~$\sigma$. We execute each configuration with the fixed seed set $\mathcal{S} = \{42, 123, 456, 789, 101112\}$ and average per-seed measurements to reduce seed-induced variance in the training features. The dataset contains 3{,}739 averaged configurations. Table~\ref{tab:data_accounting} summarizes the dataset composition.

\begin{table}[t]
\centering
\caption{Dataset summary.}
\label{tab:data_accounting}
\begin{tabular}{@{} lrr @{\hspace{2em}} lrr @{}}
\toprule
\textbf{Metric} & \textbf{Encoder} & \textbf{Decoder} & \textbf{Metric} & \textbf{Encoder} & \textbf{Decoder} \\
\midrule
Total   & 1{,}891 & 1{,}848 & Fault categories & 11 & 12 \\
Faulty  & 1{,}598 & 1{,}180 & Root causes      & 38 & 45 \\
Correct & 293     & 668     &                  &    &    \\
\midrule
\textbf{Category} & \textbf{Encoder} & \textbf{Decoder} & \textbf{Category} & \textbf{Encoder} & \textbf{Decoder} \\
\midrule
Variant    & 24  & 51  & LayerNorm & 210 & 185 \\
Score      & 278 & 142 & Embedding & 176 & 165 \\
QKV        & 247 & 184 & Output    & 192 & 180 \\
Positional & 180 & 179 & Masking   & 70  & 126 \\
FFN        & 288 & 203 & Kernel    & 96  & 63  \\
Residual   & 130 & 167 & KV Cache  & --- & 203 \\
\bottomrule
\end{tabular}
\end{table}


For \emph{fault detection} (Level~1), the label is binary. Each sample is labeled as \textit{faulty} or \textit{correct}, and corresponds to one configuration averaged across five seed runs. For \emph{fault categorization} (Level~2), the label is the injected \emph{fault category}. Encoders use 11 categories: Score, Positional, FFN, QKV, LayerNorm, Residual, Masking, Embedding, Kernel, Output, and Variant. Decoders use the same 11 plus KV Cache, covering faults in the key-value caching mechanism used during autoregressive generation. For \emph{root-cause diagnosis} (Level~3), the label is the injected \emph{root cause} within each fault category. Each category contains 2--7 root causes that distinguish structurally different faults within the same transformer component. Level~3 operates conditionally on the category predicted at Level~2. The label space is restricted to that category's root causes.

Figure~\ref{fig:defaultpp-fault-taxonomy} shows the complete fault taxonomy that defines the Level~2 and Level~3 label spaces. The seven attention-specific categories and their 25 root causes are derived from the attention fault taxonomy (Chapter~\ref{chapter_attention})~\cite{attention_taxonomy}. The five non-attention categories (Input/Embedding, Layer Normalization, FFN, Residual, and Output) and their 20 root causes are derived from prior DNN fault taxonomies~\cite{faulttaxonomy} and DeepCrime mutation operators~\cite{deepcrime}. Each root cause name identifies the fundamental mechanism or parameter that is faulty (e.g., activation function, dropout rate, weight scaling) rather than a specific corrective action. The 12 categories collectively define 38 root causes for encoders and 45 for decoders.

\begin{figure}[!htbp]
\centering
\includegraphics[width=\textwidth]{taxonomy-chapter-7-8.png}
\caption[Fault taxonomy defining the Level~2 and Level~3 label spaces.]{Complete fault taxonomy for transformer architectures. The seven attention-specific categories (bottom) and 25 root causes are from the attention fault taxonomy (Chapter~\ref{chapter_attention}). The five non-attention categories (top) and 20 root causes are from prior DNN fault taxonomies and DeepCrime.}
\label{fig:defaultpp-fault-taxonomy}
\end{figure}

The encoder and decoder subsets in Table~\ref{tab:data_accounting} are imbalanced. Encoder instances are 84.5\% faulty and decoder instances are 63.9\% faulty. We therefore treat fault detection, fault categorization, and root-cause diagnosis as imbalanced classification problems. For Levels~2 and~3, the 11--12 fault categories and their root causes also vary in support. We address this with inverse-frequency class weighting and report Macro-F1 so that low-support classes remain visible in the evaluation.

%------------------------------------------------------------------------------
\subsection{Feature Representation}
\label{sec:defaultpp_features}

Generic DNN fault features such as loss trajectories, gradient statistics, and activation counts do not distinguish transformer fault categories~\cite{frankenformer}. Corrupted QKV projections, missing causal masks, and misconfigured positional encodings leave end-task metrics undisturbed while producing characteristic patterns in attention-level and component-level measurements. We therefore define a transformer-specific feature space that targets attention distributions, projection alignments, and residual-stream behavior.

We instrument transformer training across 13 feature groups (Table~\ref{tab:defaultpp_features}). Each group measures a distinct behavioral aspect of transformer execution such as attention entropy, projection alignment, or gradient noise scale. Encoders and decoders share most groups, but decoders add cache-specific diagnostics not available in encoder instrumentation~\cite{frankenformer}. For each core metric, we extract three temporal summaries covering early training (epoch~1), mid-training (epochs~2--3), and the final epoch, with mean and slope captured for each stage. We apply coefficient-of-variation filtering (CV $< 0.01$) during preprocessing to remove near-constant columns. We organize the resulting feature vector by group for input to the group-encoding MLPs.

\begin{table}[!htbp]
\caption{Feature groups in DEFault++.}
\label{tab:defaultpp_features}
\centering
\footnotesize
\renewcommand{\arraystretch}{0.92}
\resizebox{\linewidth}{!}{%
\begin{tabular}{@{} p{3.0cm} p{3.8cm} c c p{5.5cm} @{}}
\toprule
\textbf{Feature Group} & \textbf{Core Metrics} & \textbf{Enc.} & \textbf{Dec.} & \textbf{Diagnostic Purpose} \\
\midrule
Attention statistics & Entropy, attention weight on padding tokens, head utilization (per layer) & \cmark & \cmark & Quantifies attention distribution quality. Flags collapse onto single positions, uniform dispersion, and excessive padding attention \\
QKV alignment & Q--K, Q--V, K--V cosine similarity & \cmark & \xmark & Measures projection coherence. Detects corrupted weight interactions \\
Score statistics & Mean, variance, max of raw attention logits & \cmark & \cmark & Measures attention computation stability. Sensitive to Score and Positional faults \\
Positional discrim. & Position distinguishability in representations & \cmark & \cmark & Measures whether positions remain separable. Targets Positional faults \\
FFN output magnitude & L2 norm of FFN unit outputs (per layer) & \cmark & \cmark & Measures feed-forward health independently of attention metrics \\
LayerNorm statistics & $\gamma$ norms, post-norm mean/variance & \cmark & \xmark & Reveals normalization disruptions. Primary indicator for LayerNorm faults \\
Residual stream sim. & Cosine similarity pre/post residual connection & \cmark & \xmark & Tracks skip-connection integrity. Diagnostic for Residual faults \\
Representation drift & CKA~\cite{kornblith2019similarity} layer-wise representation change & \cmark & \xmark & Measures long-range fault propagation across layers \\
Embedding statistics & Input embedding norm, token-level variance & \cmark & \xmark & Detects corruption at the model entry point. Targets Embedding faults \\
Training dynamics & Loss trajectory, gradient norm-to-noise, step time & \cmark & \cmark & Monitors optimization health. Sensitive to faults that destabilize training \\
Task metrics & Accuracy/perplexity, calibration error & \cmark & \cmark & Measures end-task impact. Separates faults with observable task-level degradation from task-invisible faults \\
Kernel timing & GPU kernel execution time, peak memory & \cmark & \cmark & Detects computational anomalies. Primary indicator for Kernel faults \\
Cache diagnostics & cache\_hidden\_sim., cache\_nll\_divergence & \xmark & \cmark & Targets KV Cache faults unique to autoregressive decoders \\
\bottomrule
\end{tabular}%
}
\end{table}

The 13 groups fall into four functional categories. \emph{Attention-level features} (attention statistics, QKV alignment, score statistics, positional discrimination) characterize the internal dynamics of multi-head attention. QKV, Score, and Positional faults share overlapping feature distributions in the attention measurement space. These three categories are the most difficult to distinguish within the Level~2 label space.

\emph{Component-level features} (FFN output magnitude, LayerNorm statistics, residual stream similarity, embedding statistics) target specific transformer units. Residual-stream cosine similarity identifies skip-connection faults. These faults perturb the residual stream at each layer they bypass. All four component-level groups are encoder-only because the decoder instrumentation in~\cite{frankenformer} exposes only layer-level aggregate outputs rather than per-component intermediate representations.

\emph{Propagation and cache features} include CKA~\cite{kornblith2019similarity} representation drift, which measures how far a fault propagates through layers, and decoder-specific cache diagnostics for KV Cache faults. We measure the cosine similarity between cached and recomputed key-value states (\texttt{cache\_hidden\_sim}) and the KL divergence between cached and fresh next-token distributions (\texttt{cache\_nll\_divergence}).

\emph{Training and system features} (loss trajectory, gradient statistics, task metrics, kernel timing) are architecture-agnostic and characterize optimization stability and computational anomalies.

% Top features per group figure removed (superseded by updated analysis).

We design the feature groups around cross-component measurements. Transformer faults propagate across heads, projections, and residual streams~\cite{hessian_analysis,roquet2024cross}, so end-task metrics alone are insufficient to localize them. We expand each core feature across training phases to preserve the temporal structure of fault development~\cite{frankenformer}.

%------------------------------------------------------------------------------
\subsection{Diagnostic Model}
\label{sec:defaultpp_reasoning}

\subsubsection{Fault Propagation Graph}
\label{sec:defaultpp_fpg}

The FPG is a deterministic directed graph that encodes how implementation faults propagate between transformer components. Prior fault diagnosis graphs are learned from data~\cite{fl4deep}. We derive the FPG from the transformer forward and backward pass without requiring training data.

We define the FPG $\mathcal{G} = (\mathcal{V}, \mathcal{E})$ where each node $v \in \mathcal{V}$ represents a transformer component. For encoders, the ten component types are embedding, positional encoding, QKV projection, score computation, attention masking, attention output, residual connections, LayerNorm, FFN, and output head. The component-level FPG instantiates 14 nodes because residual connections and LayerNorm appear independently for each sublayer. Decoders additionally include KV cache and kernel runtime, giving 15 nodes in the decoder FPG. Each directed edge $(v_i, v_j) \in \mathcal{E}$ indicates that a fault in $v_i$ propagates to $v_j$.

By analyzing the transformer's forward and backward equations, we identify eight propagation mechanisms. Let $\delta x$ denote a perturbation to variable $x$.

\begin{enumerate}[nosep]
\item \textbf{Forward sequential propagation.} If $B = f(A)$, then $\delta B = (\partial f / \partial A)\, \delta A$. Perturbation propagates along the data flow path: embedding $\to$ projections $\to$ scores $\to$ masking $\to$ weights $\to$ output $\to$ residual $\to$ LayerNorm $\to$ FFN $\to$ residual $\to$ output head.

\item \textbf{Simultaneous propagation.} QKV projections feed both score computation (via $Q, K$) and attention output (via $V$) simultaneously. A single projection fault therefore perturbs scores and values at the same time. In decoders, $K$ and $V$ additionally enter the cache.

\item \textbf{Residual bypass.} In $h = x + \text{MHA}(x)$, a perturbation $\delta x$ appears in $h$ with unit gain regardless of what the attention sublayer does. Skip connections guarantee that faults cannot be absorbed by the sublayer they bypass.

\item \textbf{Cross-layer propagation.} The residual stream propagates perturbations with at least unit gain per layer~\cite{roquet2024cross}. After $N$ layers, a fault in layer $\ell$ is present in all representations $\ell+1, \ldots, \ell+N$.

\item \textbf{Backward pass propagation.} Any fault that changes the loss $\mathcal{L}$ changes $\partial \mathcal{L} / \partial \theta_i$ for every parameter $\theta_i$, affecting weight updates of every component during training.

\item \textbf{Shared hardware propagation.} GPU kernel configuration (precision, scheduling) affects all operations that execute through that kernel. The propagation path runs through the shared execution environment, not through data flow.

\item \textbf{Architecture-wide propagation.} Some fault types inherently affect multiple components at once. A variant fault (e.g., single-head instead of multi-head attention) simultaneously changes the attention computation pattern, parameter shapes, and kernel dispatch.

\item \textbf{Cache propagation across time steps} (decoder only). The KV cache stores $K$ and $V$ across generation steps. A cache fault affects the current token and all future tokens.
\end{enumerate}

Three nonlinear operations limit propagation at specific edges. Softmax maps perturbations into $[0,1]$ with a sum-to-one constraint, which reduces large-magnitude deviations. LayerNorm~\cite{ba2016layer} rescales outputs to zero mean and unit variance, which removes additive shifts. Activation functions (ReLU, GELU) suppress perturbations in their saturation regions and pass them in their linear regions. We annotate these limiting operations on the corresponding edges of the FPG.

At the component level, the encoder FPG contains 14 nodes and 31 edges (one node per component listed above, with edges derived from the eight mechanisms). The decoder FPG contains 15 nodes (adding KV cache) and 34 edges. For the neural model, we collapse the component-level FPG to feature-group level by mapping each component to its feature group (Section~\ref{sec:defaultpp_features}). This produces a group-level adjacency matrix with self-loops. Each component type maps to the feature group that instruments it. For example, QKV projection maps to the QKV alignment group, score computation maps to the score statistics group, and so on following Table~\ref{tab:defaultpp_features}. Figure~\ref{fig:defaultpp_fpg} shows the resulting graph.

\begin{figure}[!htbp]
\centering
\includegraphics[width=0.85\linewidth]{fig_fpg.pdf}
\caption{The Fault Propagation Graph (FPG) for transformer architectures.}
\label{fig:defaultpp_fpg}
\end{figure}

\subsubsection{Feature-Group Encoding}
\label{sec:defaultpp_encoding}

We encode each feature group through a dedicated MLP. We keep groups structurally separated so that message passing on the FPG can propagate information between them and so that the explanation step can examine each group's embedding independently. We use trainable MLPs rather than a flat classifier because the prototype-distance objective and the contrastive separation loss both require gradients to flow through the group-level weights during training.

Given the feature vector $\mathbf{z} \in \mathbb{R}^d$ partitioned into $G$ groups $\{\mathbf{z}_{g_1}, \ldots, \mathbf{z}_{g_G}\}$ corresponding to the feature groups in Table~\ref{tab:defaultpp_features}, each group $g$ is encoded by a dedicated MLP (Equation~\ref{eq:group_encoding}).
\begin{equation}
\mathbf{h}_g = \text{MLP}_g(\mathbf{z}_g) \in \mathbb{R}^{h}, \quad g = 1, \ldots, G,
\label{eq:group_encoding}
\end{equation}
where $h = 32$ is the hidden dimension per group. The group embeddings form the matrix $\mathbf{H} = [\mathbf{h}_1; \ldots; \mathbf{h}_G] \in \mathbb{R}^{G \times h}$.

A projection layer maps the concatenated group embeddings to a shared representation (Equation~\ref{eq:projection}).
\begin{equation}
\mathbf{z}_{\text{proj}} = W_{\text{proj}} \, \text{vec}(\mathbf{H}) \in \mathbb{R}^{e},
\label{eq:projection}
\end{equation}
where $e = 64$ is the embedding dimension and $\text{vec}(\mathbf{H}) \in \mathbb{R}^{Gh}$ denotes the row-major vectorization. Levels~1 and~2 use the pooled projection $\mathbf{z}_{\text{proj}}$ for detection and categorization. Level~3 uses the group embedding matrix $\mathbf{H}$ directly to preserve group-level structure for root-cause separation and distance-based explanation.

\subsubsection{Message Passing on the FPG}
\label{sec:defaultpp_message_passing}

We update each group embedding by aggregating information from the groups connected via FPG edges.

Given the group-level adjacency matrix $\mathbf{A} \in \mathbb{R}^{G \times G}$ derived from the FPG (Section~\ref{sec:defaultpp_fpg}), we compute updated group embeddings following graph convolutional networks~\cite{kipf2017semi}.
\begin{equation}
\mathbf{H}' = \text{ReLU}\!\left(\hat{\mathbf{A}}\, \mathbf{H}\, W_{\text{msg}}\right),
\label{eq:message_passing}
\end{equation}
where $\hat{\mathbf{A}}$ is the row-normalized adjacency matrix (with self-loops) and $W_{\text{msg}} \in \mathbb{R}^{h \times h}$. We use row normalization rather than symmetric normalization because feature groups differ in the number of FPG propagation connections. Row normalization prevents high-degree groups from over-averaging their neighbors. The updated group embeddings $\mathbf{H}'$ replace the pre-message-passing $\mathbf{H}$ in both the projection and the Level~3 prototype matching. We disable message passing in the ablation that isolates the contribution of FPG structure (RQ$_2$). We apply three rounds of message passing, each with a dedicated learnable message function. We selected the number of rounds via grid search on the inner validation fold (Appendix~\ref{app:defaultpp-parameters}).

For example, a QKV projection fault propagates to attention scores and the residual stream according to Mechanisms~1 and~3. Message passing along these edges aggregates that multi-group information into the QKV group embedding before classification.

\subsubsection{Hierarchical Classification}
\label{sec:defaultpp_stages}

All three classification levels share the same feature-encoding pathway (group encoding $+$ optional message passing on the FPG $+$ projection). Each level adds a task-specific head.

\paragraph{Level~1: Fault detection.}
A two-layer MLP maps $\mathbf{z}_{\text{proj}}$ to binary logits (faulty vs.\ correct).

\paragraph{Level~2: Fault categorization.}
A two-layer MLP maps $\mathbf{z}_{\text{proj}}$ to $C$-class logits, where $C = 11$ for encoders and $C = 12$ for decoders. During training, the class-weighted cross-entropy is computed only over samples labeled faulty.

\paragraph{Level~3: Root-cause diagnosis.}
Level~3 distinguishes 2--5 root causes within a fault category on limited per-class training samples. A standard softmax classifier risks overfitting to majority root causes in this regime. We therefore combine a cross-entropy head with a prototypical classifier~\cite{snell2017prototypical,kamp1995prototype}. The prototypical classifier reduces the per-class parameter burden by comparing each sample to class-mean embeddings rather than learning a separate weight vector per class.

First, a per-category classification head maps $\mathbf{z}_{\text{proj}}$ to root-cause logits within each category. During training, the cross-entropy loss is computed only over samples belonging to the true category. Second, a prototypical classifier~\cite{snell2017prototypical,lin2022prototypical} operates in the embedding space $\mathbf{H} \in \mathbb{R}^{G \times h}$. Let $\mathcal{D}_{c,r}$ denote the set of training samples with category $c$ and root cause $r$. For each root cause $r$ within category $c$, we compute a prototype $\symbf{\pi}_{c,r}$ as the mean group embedding (Equation~\ref{eq:prototype}).
\begin{equation}
\symbf{\pi}_{c,r} = \frac{1}{|\mathcal{D}_{c,r}|} \sum_{i \in \mathcal{D}_{c,r}} \mathbf{H}_i \in \mathbb{R}^{G \times h}.
\label{eq:prototype}
\end{equation}

The distance from a sample to a prototype decomposes additively by group (Equation~\ref{eq:proto_distance}).
\begin{equation}
d(\mathbf{H}, \symbf{\pi}_{c,r}) = \sum_{g=1}^{G} \| \mathbf{h}_g - \symbf{\pi}_{c,r,g} \|^2 = \sum_{g=1}^{G} d_g,
\label{eq:proto_distance}
\end{equation}
where $d_g = \| \mathbf{h}_g - \symbf{\pi}_{c,r,g} \|^2$ is the distance contribution of group $g$, which forms the basis for the FPG-based explanation described in Section~\ref{sec:defaultpp_explainability}.

At inference time, Level~3 is invoked only after Level~2 predicts a fault category. Root-cause diagnosis is then restricted to that predicted category. The predicted root cause is the one whose prototype is nearest in the grouped space. The per-category cross-entropy heads and the prototype-matching loss $\mathcal{L}_{\text{pm}}$ both train on Level~3 samples. The separation loss $\mathcal{L}_{\text{sep}}$ provides additional within-category structure (Section~\ref{sec:defaultpp_training}).

\subsubsection{Training Objective}
\label{sec:defaultpp_training}

The entire network, including the shared feature-group MLPs, is trained end-to-end. Gradients from all three levels backpropagate through the shared encoding pathway. The Level~3 losses therefore refine the shared feature embeddings alongside the detection and categorization objectives.

Each level uses a dedicated loss. Level~1 uses binary cross-entropy $\mathcal{L}_{\text{detect}}$ over all samples. Level~2 uses class-weighted cross-entropy $\mathcal{L}_{\text{cat}}$ over fault categories, applied to faulty samples only. Level~3 uses two losses jointly: per-category cross-entropy $\mathcal{L}_{\text{rc}}$ over root causes, and a root-cause separation loss $\mathcal{L}_{\text{sep}}$ that enforces within-category discrimination in the embedding space.

Cross-entropy trains per-category heads but does not explicitly structure the embedding space for within-category separation. We therefore define $\mathcal{L}_{\text{sep}}$ with two components. The first is a contrastive term that operates on the flattened group embeddings $\text{vec}(\mathbf{H})$. For each sample $i$ with category label $c_i$ and root-cause label $r_i$, positive pairs share both $(c_i, r_i)$, while negative pairs share $c_i$ but differ in root cause. Following supervised contrastive learning~\cite{khosla2020supervised}, the contrastive term within category $c$ is
\begin{equation}
\mathcal{L}_{\text{ctr}}^{(c)} = -\frac{1}{|\mathcal{P}_c|} \sum_{(i,j) \in \mathcal{P}_c} \log \frac{\exp(\text{sim}(\mathbf{h}_i, \mathbf{h}_j) / \tau)}{\sum_{k : c_k = c_i} \exp(\text{sim}(\mathbf{h}_i, \mathbf{h}_k) / \tau)},
\label{eq:contrastive}
\end{equation}
where $\mathcal{P}_c$ is the set of positive pairs within category $c$, $\text{sim}(\cdot, \cdot)$ is cosine similarity on the normalized flattened group embeddings, and $\tau = 0.1$ is the temperature. We use cosine similarity here because the contrastive objective targets directional alignment of embeddings, not their absolute magnitude~\cite{wang2021understanding,islam2023self}. The prototype classifier in Equation~\ref{eq:proto_distance} uses squared Euclidean distance instead, because prototype matching requires spatial proximity in the grouped embedding space rather than directional agreement.

We define the joint training objective $\mathcal{L}_{\text{DEFault++}}$ in Equation~\ref{eq:total_loss}.
\begin{equation}
\mathcal{L}_{\text{DEFault++}} = \mathcal{L}_{\text{detect}} + \alpha \, \mathcal{L}_{\text{cat}} + \lambda \, \mathcal{L}_{\text{rc}} + \mathcal{L}_{\text{sep}},
\label{eq:total_loss}
\end{equation}
where $\alpha$ and $\lambda$ are level-specific weighting coefficients. The root-cause separation loss $\mathcal{L}_{\text{sep}}$ is applied only to faulty samples within the predicted fault category and is defined as
\begin{equation}
\mathcal{L}_{\text{sep}} = \beta \, \mathcal{L}_{\text{ctr}} + \gamma \, \mathcal{L}_{\text{pm}},
\label{eq:separation_loss}
\end{equation}
where $\mathcal{L}_{\text{ctr}}$ is the contrastive term from Equation~\ref{eq:contrastive}, and $\mathcal{L}_{\text{pm}}$ is a prototype-matching loss defined as
\begin{equation}
\mathcal{L}_{\text{pm}} = -\frac{1}{|\mathcal{B}|} \sum_{i \in \mathcal{B}} \log \frac{\exp(-d(\mathbf{H}_i, \symbf{\pi}_{c_i, r_i}) / \tau_p)}{\sum_{r' \in \mathcal{R}_{c_i}} \exp(-d(\mathbf{H}_i, \symbf{\pi}_{c_i, r'}) / \tau_p)},
\label{eq:proto_matching}
\end{equation}
where $\mathcal{B}$ is the training batch, $d(\cdot, \cdot)$ is the grouped squared distance from Equation~\ref{eq:proto_distance}, $\mathcal{R}_{c_i}$ is the set of root causes within category $c_i$, and $\tau_p$ is a temperature parameter. By converting negative prototype distances into class logits, $\mathcal{L}_{\text{pm}}$ trains $\mathbf{H}$ toward the distance-based explanation in Equation~\ref{eq:proto_distance}.

We set $\alpha = 1.0$, $\lambda = 1.0$, $\beta = 0.5$, and $\gamma = 0.3$, where $\beta$ weights the contrastive term $\mathcal{L}_{\text{ctr}}$ and $\gamma$ weights the prototype-matching term $\mathcal{L}_{\text{pm}}$ in $\mathcal{L}_{\text{sep}}$. We selected all four coefficients via grid search on the inner validation fold (Appendix~\ref{app:defaultpp-parameters}). Training details appear in Section~\ref{sec:defaultpp_experiment}.

\subsubsection{FPG-Based Explanation}
\label{sec:defaultpp_explainability}

For each diagnosed sample, the distance decomposition in Equation~\ref{eq:proto_distance} produces a vector $[d_1, \ldots, d_G]$ where $d_g$ is group $g$'s distance contribution to the predicted root-cause prototype. We normalize these distances to obtain a fractional salience score per component
\begin{equation}
w_g = \frac{d_g}{\sum_{g'=1}^{G} d_{g'}}, \quad g = 1, \ldots, G.
\label{eq:explanation}
\end{equation}

The groups are aligned with FPG nodes, so each $w_g$ maps directly to a transformer subsystem. Because $w_g$ is computed from distances in the post-message-passing embedding space, it measures fault effects as they propagate across components rather than deviations in raw feature values at a single source. A large $w_g$ indicates that group $g$ is the subsystem most affected by the fault after propagation~\cite{wolf2024keep}. The group ablation in RQ$_3$ (Section~\ref{sec:defaultpp_rq3}) tests this interpretation by zeroing the highest-$w_g$ group embeddings and measuring the resulting drop in root-cause diagnosis performance.

%==============================================================================
\section{Experiment}
\label{sec:defaultpp_experiment}

We evaluate DEFault++ on the labeled dataset through three research questions. Encoder and decoder architectures are evaluated separately because their feature and label spaces differ.

\begin{rqlist}
\item \textbf{RQ$\mathbf{_1}$:} How does DEFault++ perform in detecting and categorizing transformer faults?
\item \textbf{RQ$\mathbf{_2}$:} How do message passing on the FPG and the separation loss contribute to diagnostic performance?
\item \textbf{RQ$\mathbf{_3}$:} How does DEFault++ perform in root-cause diagnosis, and how faithful are the explanations?
\end{rqlist}

%------------------------------------------------------------------------------
\subsection{Evaluation Criteria}
\label{sec:defaultpp_cv}

We use nested grouped cross-validation with $k=5$ outer folds, grouping instances by $(\text{model}, \text{dataset}, \text{seed})$ to prevent information leakage from correlated traces. The inner split uses \texttt{StratifiedGroupKFold} for model selection. We restrict all preprocessing (e.g., scaling, thresholding, and standardization) and hyperparameter grid search strictly to the inner training fold. Finally, we fix the outer fold assignments across all evaluated methods to ensure each method is evaluated on the same held-out samples. \footnote{Replication package: \url{https://github.com/SigmaJahan/DEFaultplusplus-Transformer-Debugging}}

We evaluate DEFault++ using four standard metrics. Let $\text{TP}$, $\text{FP}$, $\text{TN}$, and $\text{FN}$ denote the true positive, false positive, true negative, and false negative counts, respectively. Precision ($P$) and recall ($R$) are defined as
\begin{equation}
P = \frac{\text{TP}}{\text{TP} + \text{FP}}, \qquad R = \frac{\text{TP}}{\text{TP} + \text{FN}}.
\label{eq:prec_recall}
\end{equation}
The F1-score is the harmonic mean of precision and recall:
\begin{equation}
\text{F1} = \frac{2 \cdot P \cdot R}{P + R}.
\label{eq:f1}
\end{equation}
The class distributions are imbalanced at all three levels. We therefore report Macro-F1~\cite{naidu2023review}, which averages F1 across classes:
\begin{equation}
\text{Macro-F1} = \frac{1}{C} \sum_{c=1}^{C} \text{F1}_c,
\label{eq:macro_f1}
\end{equation}
where $C$ is the number of classes and $\text{F1}_c$ is the F1-score for class $c$. For fault detection (Level~1), we additionally report AUROC~\cite{bradley1997use}, which measures discriminative ability across all classification thresholds.

A true positive for fault detection means that a faulty configuration is correctly identified, regardless of its fault category. A true positive for fault categorization means that the category of the fault is identified correctly, regardless of the root cause. A true positive for root-cause diagnosis requires both the Level~2 category and the Level~3 root cause to be identified correctly in the same inference path.

\subsection{Implementation}

We implement the hierarchical model in PyTorch. The encoder uses separate MLPs for each feature group with hidden dimension $h = 32$ and a shared projection to $e = 64$ dimensions. We apply three rounds of message passing on the FPG with the group-level adjacency matrix. Training uses Adam~\cite{kingma2015adam} with learning rate $10^{-3}$, weight decay $10^{-4}$, and batch size 256. We train for up to 150 epochs with early stopping on the validation split. The loss weights are $\alpha = 1.0$, $\lambda = 1.0$, $\beta = 0.5$, and $\gamma = 0.3$, with temperature $\tau = 0.1$. We address fault-category and root-cause imbalance through inverse-frequency weighting in the corresponding classification losses and fix all random seeds for reproducibility. Appendix~\ref{app:defaultpp-parameters} (Table~\ref{tab:defaultpp-hyperparameters}) gives the full hyperparameter configuration. The replication package~\cite{defaultpp_repo} contains the implementation, hyperparameter files, and trained models.

\subsection{Ablation Design}

To isolate the contribution of each architectural component, we evaluate ablation variants on the same folds. We ablate two components: message passing on the FPG and the separation loss $\mathcal{L}_{\text{sep}}$. The separation loss is designed for within-category root-cause discrimination and applies only at Level~3. We therefore define two variants for Levels~1 and~2 and four variants for Level~3:

\begin{enumerate}[nosep]
\item \textbf{Base}. Feature-group MLPs with projection, no message passing on the FPG, and no separation loss ($\beta = 0$, $\gamma = 0$). This variant tests the hierarchical classifier alone.
\item \textbf{+FPG}. Base plus message passing on the FPG. This isolates the effect of propagation-informed feature interaction. At Levels~1 and~2, +FPG is the complete model because $\mathcal{L}_{\text{sep}}$ does not apply to detection or categorization.
\item \textbf{+Sep (Level~3 only)}. Base plus the separation loss ($\beta = 0.5$, $\gamma = 0.3$), without message passing on the FPG. This isolates the effect of within-category separation on root-cause diagnosis.
\item \textbf{Full (FPG+Sep)}. Message passing on the FPG together with the separation loss at Level~3. This is the complete DEFault++ model.
\end{enumerate}

\subsection{Baseline Approaches}

We compare DEFault++ against four prior techniques only on the directly comparable Level~1 detection task. For each outer fold, the baselines are evaluated on the same train/test partition as DEFault++. Rule-based techniques are executed with their published rule definitions and thresholds after mapping the required observables to the closest available trace features in our dataset. Learned baselines are trained as binary detectors on the Level~1 labels of each outer-training fold using the model families provided in their replication packages~\cite{autotrainer,deepdiagnosis,deepfd,default}. We do not recalibrate published rule thresholds on our dataset, and we do not claim Level~2 or Level~3 comparability when the native label space does not map validly to our transformer component categories or root causes.

\begin{enumerate}[label=(\roman*),nosep]
\item \textbf{AutoTrainer}~\cite{autotrainer} monitors five training symptoms (vanishing gradient, exploding gradient, dying ReLU, oscillating loss, slow convergence) via threshold rules applied during training. We map each rule to the corresponding features in our dataset and apply the published thresholds without recalibration. AutoTrainer produces binary detection decisions but its symptom types do not map to our fault categories.

\item \textbf{DeepDiagnosis}~\cite{deepdiagnosis} extends AutoTrainer to eight symptom rules (adding unchanged weights, saturated activation, and out-of-range output) with a decision tree that maps detected symptoms to seven fix types. We apply the same feature-mapping strategy and thresholds. Its symptom labels likewise do not map to our fault categories.

\item \textbf{DeepFD}~\cite{deepfd} combines KNN, Decision Tree, and Random Forest classifiers with union voting to diagnose faults from runtime features. Its five diagnosis classes (e.g., incorrect learning rate, wrong activation function) do not correspond to our fault categories. We evaluate DeepFD for binary detection only.

\item \textbf{DEFault}~\cite{default} (Chapter~\ref{chapter_default}) uses a hierarchical Random Forest classifier that combines static and dynamic features to detect and categorize faults across seven coarse DNN-level categories (learning rate, activation function, loss function, optimizer, layer configuration, regularization, data processing). These categories operate at a different abstraction level. We evaluate DEFault for binary detection only.
\end{enumerate}

Table~\ref{tab:baseline_coverage} summarizes the coverage of each technique across our diagnostic stages.

\begin{table}[!htbp]
\caption{Baseline coverage across DEFault++ diagnostic stages.}
\label{tab:baseline_coverage}
\centering
\footnotesize
\setlength{\tabcolsep}{4pt}
\begin{tabular}{@{} l l c c c @{}}
\toprule
\textbf{Technique} & \textbf{Paradigm} & \textbf{Detection} & \textbf{Categorization} & \textbf{Root-Cause} \\
\midrule
AutoTrainer~\cite{autotrainer}     & Rule-based (5 rules)      & \cmark & N/A & N/A \\
DeepDiagnosis~\cite{deepdiagnosis} & Rule-based (8 rules)      & \cmark & N/A & N/A \\
DeepFD~\cite{deepfd}               & Flat ensemble (KNN+DT+RF) & \cmark & N/A & N/A \\
DEFault~\cite{default}             & Hierarchical RF            & \cmark & N/A & N/A \\
\midrule
\textbf{DEFault++}          & Neural hierarchical        & \cmark & \cmark & \cmark \\
\bottomrule
\end{tabular}
\end{table}

%------------------------------------------------------------------------------
\subsection{Evaluation of DEFault++}

\subsubsection{\textbf{RQ$\mathbf{_1}$:} Fault Detection and Fault Categorization}
\label{sec:defaultpp_rq1}

We evaluate DEFault++ on fault detection and fault categorization. Tables~\ref{tab:perf_level1} and~\ref{tab:perf_level2_overall} report the results.

\begin{table}[!htbp]
\centering
\caption{DEFault++ performance on fault detection.}
\label{tab:perf_level1}
\begin{tabular}{l c c}
\toprule
\textbf{Metric} & \textbf{Encoder} & \textbf{Decoder} \\
\midrule
AUROC & 0.9660 $\pm$ 0.0044 & 0.9620 $\pm$ 0.0049 \\
Accuracy & 0.9455 $\pm$ 0.0060 & 0.9107 $\pm$ 0.0041 \\
Macro-F1 & 0.8238 $\pm$ 0.0147 & 0.9088 $\pm$ 0.0041 \\
Weighted-F1 & 0.9477 $\pm$ 0.0051 & 0.9104 $\pm$ 0.0040 \\
Precision & 0.8014 $\pm$ 0.0210 & 0.9125 $\pm$ 0.0049 \\
Recall & 0.8520 $\pm$ 0.0089 & 0.9064 $\pm$ 0.0036 \\
\bottomrule
\end{tabular}
\end{table}

\begin{table}[!htbp]
\centering
\caption{DEFault++ performance on fault categorization.}
\label{tab:perf_level2_overall}
\begin{tabular}{l c c}
\toprule
\textbf{Metric} & \textbf{Encoder} & \textbf{Decoder} \\
\midrule
Accuracy & 0.8522 $\pm$ 0.0126 & 0.8842 $\pm$ 0.0094 \\
Macro-F1 & 0.8497 $\pm$ 0.0129 & 0.8679 $\pm$ 0.0092 \\
Weighted-F1 & 0.8531 $\pm$ 0.0119 & 0.8850 $\pm$ 0.0091 \\
Precision & 0.8662 $\pm$ 0.0126 & 0.8794 $\pm$ 0.0080 \\
Recall & 0.8403 $\pm$ 0.0148 & 0.8649 $\pm$ 0.0097 \\
\bottomrule
\end{tabular}
\end{table}

We observe AUROC of 0.966 on encoders and 0.962 on decoders for fault detection (Table~\ref{tab:perf_level1}), and Macro-F1 of 0.850 and 0.868 for fault categorization (Table~\ref{tab:perf_level2_overall}). Detection recall of 0.852 (encoder) and 0.906 (decoder) indicates that DEFault++ identifies most faulty configurations before they reach Level~2. We also find that categorization precision and recall remain balanced (0.866/0.840 for encoders, 0.879/0.865 for decoders), which suggests that the model does not sacrifice one for the other. Table~\ref{tab:rq1_perclass} breaks down categorization performance by fault category.

\begin{table}[!htbp]
\centering
\caption{Per-category categorization performance for DEFault++.}
\label{tab:rq1_perclass}
\small
\begin{tabular}{l c c c c c c}
\toprule
& \multicolumn{3}{c}{\textbf{Encoder}} & \multicolumn{3}{c}{\textbf{Decoder}} \\
\cmidrule(lr){2-4} \cmidrule(lr){5-7}
\textbf{Category} & \textbf{F1} & \textbf{Prec.} & \textbf{Rec.} & \textbf{F1} & \textbf{Prec.} & \textbf{Rec.} \\
\midrule
Embedding & 0.838 & 0.852 & 0.829 & 0.862 & 0.833 & 0.897 \\
FFN & 0.836 & 0.807 & 0.871 & 0.718 & 0.673 & 0.777 \\
Kernel & 0.771 & 0.814 & 0.740 & 0.765 & 0.848 & 0.700 \\
KV Cache & --- & --- & --- & 0.758 & 0.716 & 0.822 \\
LayerNorm & 0.790 & 0.775 & 0.820 & 0.843 & 0.937 & 0.766 \\
Masking & 0.917 & 0.952 & 0.886 & 0.890 & 0.927 & 0.857 \\
Output & 0.889 & 0.927 & 0.864 & 0.975 & 0.980 & 0.969 \\
Positional & 0.919 & 0.950 & 0.891 & 0.927 & 0.949 & 0.911 \\
QKV & 0.825 & 0.863 & 0.792 & 0.789 & 0.772 & 0.815 \\
Residual & 0.961 & 0.965 & 0.958 & 0.995 & 0.995 & 0.995 \\
Score & 0.861 & 0.847 & 0.876 & 0.946 & 0.937 & 0.957 \\
Variant & 0.741 & 0.776 & 0.717 & 0.948 & 0.986 & 0.914 \\
\midrule
\textbf{Macro avg.} & 0.850 & 0.866 & 0.840 & 0.868 & 0.879 & 0.865 \\
\bottomrule
\end{tabular}
\end{table}

We identify Residual and Positional as the strongest categories. Residual achieves near-perfect F1 on both architectures (0.961 encoder, 0.995 decoder), which is consistent with the distinctive skip-connection signature captured by the residual-stream feature group. On encoders, Variant (0.741) and Kernel (0.771) report the lowest F1. On decoders, FFN (0.718) and KV Cache (0.758) are the most difficult. We also observe that Variant differs sharply across architectures (F1 0.741 encoder vs.\ 0.948 decoder). One plausible explanation is that encoder Variant faults have only 24 instances (Table~\ref{tab:data_accounting}), the smallest category, while decoder Variant faults have 51 instances and benefit from stronger cache-related discriminative features. Figure~\ref{fig:defaultpp_roc} shows the detection ROC curves.

\begin{figure}[!htbp]
\centering
\includegraphics[width=0.9\linewidth]{fig_RQ1_roc_detection.png}
\caption{ROC curves for fault detection on encoder and decoder architectures.}
\label{fig:defaultpp_roc}
\end{figure}

Figures~\ref{fig:defaultpp_training_encoder} and~\ref{fig:defaultpp_training_decoder} show that training loss converges by ${\sim}$epoch 100, at which point validation Detection F1 and Category F1 also stabilize. This early convergence suggests that the learned feature space is expressive enough for the three-level hierarchy without extended training.

\begin{figure}[!htbp]
\centering
\includegraphics[width=0.9\linewidth]{fig_training_loss_encoder.png}
\caption{Training dynamics for DEFault++ on the encoder architecture.}
\label{fig:defaultpp_training_encoder}
\end{figure}

\begin{figure}[!htbp]
\centering
\includegraphics[width=0.9\linewidth]{fig_training_loss_decoder.png}
\caption{Training dynamics for DEFault++ on the decoder architecture.}
\label{fig:defaultpp_training_decoder}
\end{figure}

\FloatBarrier

\paragraph{Baseline comparison.}
We compare DEFault++ against four baselines on the binary detection task (Table~\ref{tab:rq1_detection_comparison}).
DEFault++ achieves the highest AUROC on both architectures (0.9660 encoder, 0.9620 decoder).
DEFault scores 0.8610/0.9030 and DeepFD 0.8340/0.8810, placing both above the threshold-based methods.
Macro-F1 follows the same ranking.
DEFault++ reaches 0.8238/0.9088, compared with 0.5660/0.7020 for DEFault and 0.5140/0.6550 for DeepFD.

We find that DEFault and DeepFD can still separate faulty from correct configurations, as binary detection does not require alignment with our 11/12 transformer component categories.
We attribute the gap relative to DEFault++ to the feature space each method uses.
DEFault relies on a hierarchical random forest over general DNN behavioral features, which captures broad fault signatures but does not model transformer-specific structures such as attention distributions or projection alignments.
DeepFD uses a flat ensemble over DNN-level fault categories.
This design does not provide the resolution to distinguish faults that manifest in transformer-specific components.

DeepDiagnosis and AutoTrainer produce the lowest Macro-F1 scores (0.3380/0.4820 and 0.2840/0.4270 respectively).
Both methods apply fixed threshold rules over generic DNN training symptoms.
DeepDiagnosis encodes eight such rules and AutoTrainer five, neither of which targets the attention, projection, or cache behaviors that characterize faults in transformer architectures.
We find that this design scope accounts for the larger gap relative to DEFault++.

\begin{table*}[!htbp]
\centering
\caption{Fault-detection comparison between DEFault++ and baselines.}
\label{tab:rq1_detection_comparison}
\footnotesize
\begin{adjustbox}{max width=\textwidth}
\begin{tabular}{l l c c c c c c}
\toprule
\textbf{Method} & \textbf{Type} & \textbf{AUROC (Enc)} & \textbf{M-F1 (Enc)} & \textbf{Acc.\ (Enc)} & \textbf{AUROC (Dec)} & \textbf{M-F1 (Dec)} & \textbf{Acc.\ (Dec)} \\
\midrule
\textbf{DEFault++} & Neural hierarchical & $\mathbf{0.9660}_{\pm .004}$ & $\mathbf{0.8238}_{\pm .015}$ & $\mathbf{0.9455}_{\pm .006}$ & $\mathbf{0.9620}_{\pm .005}$ & $\mathbf{0.9088}_{\pm .004}$ & $\mathbf{0.9107}_{\pm .004}$ \\
DEFault~\cite{default} & Hierarchical RF & $0.8610_{\pm .018}$ & $0.5660_{\pm .021}$ & $0.8120_{\pm .013}$ & $0.9030_{\pm .012}$ & $0.7020_{\pm .016}$ & $0.8460_{\pm .010}$ \\
DeepFD~\cite{deepfd} & Flat ensemble & $0.8340_{\pm .021}$ & $0.5140_{\pm .024}$ & $0.7810_{\pm .015}$ & $0.8810_{\pm .014}$ & $0.6550_{\pm .018}$ & $0.8190_{\pm .011}$ \\
DeepDiagnosis~\cite{deepdiagnosis} & Rule-based (8 rules) & --- & $0.3380_{\pm .028}$ & $0.7240_{\pm .019}$ & --- & $0.4820_{\pm .022}$ & $0.7710_{\pm .016}$ \\
AutoTrainer~\cite{autotrainer} & Rule-based (5 rules) & --- & $0.2840_{\pm .031}$ & $0.6970_{\pm .021}$ & --- & $0.4270_{\pm .025}$ & $0.7440_{\pm .018}$ \\
\bottomrule
\end{tabular}
\end{adjustbox}

\smallskip
\noindent{\footnotesize Rule-based techniques produce binary decisions without probabilistic scores. AUROC is therefore not applicable (---). All values are mean $\pm$ std over five outer folds. Learned methods are fit on the training portion of each outer fold, and rule-based methods are applied directly to the corresponding held-out fold.}
\end{table*}

We do not report Level~2 or Level~3 comparisons against these baselines.
AutoTrainer and DeepDiagnosis generate generic symptom labels that do not map to any fault category in our taxonomy.
DeepFD and DEFault use DNN-level fault categories that were defined independently of transformer component boundaries.
A numeric comparison at those levels would equate non-equivalent labels, thus we report baseline coverage in Table~\ref{tab:baseline_coverage} instead.

\begin{rqsummary}
\textbf{Summary of RQ$\mathbf{_1}$:} DEFault++ outperforms existing fault detection techniques, achieving an AUROC above 0.96 and a Macro-F1 above 0.85 across both architectures. Residual and Positional categories score highest due to distinct feature signatures, while Variant, Kernel, and FFN score lower due to limited support and overlapping distributions.
\end{rqsummary}

%------
\subsubsection{\textbf{RQ$\mathbf{_2}$:} Ablation Analysis}
\label{sec:defaultpp_rq2}

We examine the contribution of message passing on the FPG and the separation loss through ablation. Tables~\ref{tab:abl_l1} and~\ref{tab:abl_l2} compare Base and +FPG at Levels~1 and~2, where the separation loss does not apply. Table~\ref{tab:abl_l3} compares all four variants at Level~3.

\begin{table}[t]
\centering
\caption{Ablation study of Level~1: Fault Detection.}
\label{tab:abl_l1}
\small
\begin{tabular}{l l c c c c c c}
\toprule
\textbf{Arch.} & \textbf{Variant} & \textbf{AUROC} & \textbf{Acc.} & \textbf{F1} & \textbf{F1$_w$} & \textbf{Prec.} & \textbf{Rec.} \\
\midrule
Encoder & Base & 0.921 & 0.922 & 0.782 & 0.924 & 0.752 & 0.816 \\
        & + FPG & 0.966 & 0.946 & 0.824 & 0.948 & 0.801 & 0.852 \\
\midrule
Decoder & Base & 0.949 & 0.892 & 0.890 & 0.892 & 0.890 & 0.891 \\
        & + FPG & 0.962 & 0.911 & 0.909 & 0.910 & 0.912 & 0.906 \\
\bottomrule
\end{tabular}
\end{table}

\begin{table}[t]
\centering
\caption{Ablation study of Level~2: Fault Categorization.}
\label{tab:abl_l2}
\small
\begin{tabular}{l l c c c c c}
\toprule
\textbf{Arch.} & \textbf{Variant} & \textbf{Acc.} & \textbf{F1} & \textbf{F1$_w$} & \textbf{Prec.} & \textbf{Rec.} \\
\midrule
Encoder & Base & 0.743 & 0.742 & 0.745 & 0.757 & 0.731 \\
        & + FPG & 0.852 & 0.850 & 0.853 & 0.866 & 0.840 \\
\midrule
Decoder & Base & 0.790 & 0.788 & 0.791 & 0.797 & 0.783 \\
        & + FPG & 0.884 & 0.868 & 0.885 & 0.879 & 0.865 \\
\bottomrule
\end{tabular}
\end{table}

\begin{table}[t]
\centering
\caption{Ablation study of Level~3: Root Cause Diagnosis.}
\label{tab:abl_l3}
\small
\begin{tabular}{l l c c c c c}
\toprule
\textbf{Arch.} & \textbf{Variant} & \textbf{Acc.} & \textbf{F1} & \textbf{F1$_w$} & \textbf{Prec.} & \textbf{Rec.} \\
\midrule
Encoder & Base & 0.702 & 0.701 & 0.711 & 0.883 & 0.690 \\
        & + FPG & 0.752 & 0.756 & 0.763 & 0.910 & 0.752 \\
        & + Seperation loss & 0.778 & 0.812 & 0.819 & 0.933 & 0.776 \\
        &  FPG+Sep & 0.795 & 0.851 & 0.857 & 0.941 & 0.793 \\
\midrule
Decoder & Base & 0.735 & 0.734 & 0.742 & 0.896 & 0.721 \\
        & + FPG & 0.781 & 0.781 & 0.788 & 0.925 & 0.781 \\
        & + Seperation loss & 0.810 & 0.833 & 0.840 & 0.940 & 0.809 \\
        & FPG+Sep & 0.826 & 0.867 & 0.875 & 0.948 & 0.823 \\
\bottomrule
\end{tabular}
\end{table}

Message passing on the FPG improves both detection and categorization. At Level~1, it raises encoder Macro-F1 from 0.782 to 0.824 (+0.042) and decoder Macro-F1 from 0.890 to 0.909 (+0.019). At Level~2, the gains are larger: encoder Macro-F1 rises from 0.742 to 0.850 (+0.108) and decoder Macro-F1 from 0.788 to 0.868 (+0.080). FPG routes fault information across feature groups during message aggregation, producing more discriminative representations for detection and categorization.

At Level~3, where root causes within the same fault category must be separated, the separation loss becomes the dominant contributor. Without FPG, the separation loss alone lifts encoder Macro-F1 from 0.701 to 0.812 (+0.111) and decoder Macro-F1 from 0.734 to 0.833 (+0.099). FPG alone reaches 0.756 and 0.781. The full model combines both components and achieves the best Level~3 results on both architectures, with encoder Macro-F1 of 0.851 and decoder Macro-F1 of 0.867. FPG supplies cross-group feature interactions through message propagation, and the separation loss enforces the fine-grained boundaries that root-cause discrimination requires.

\FloatBarrier

\begin{rqsummary}
\textbf{Summary of RQ$\mathbf{_2}$:} Message passing on the FPG improves all three levels, with the largest gains at Level~2 categorization. The separation loss, applied only at Level~3, is the dominant contributor to root-cause discrimination. Their combination produces the best root-cause diagnosis results on both architectures.
\end{rqsummary}

%------
\subsubsection{\textbf{RQ$\mathbf{_3}$:} Root-Cause Diagnosis and Explanation}
\label{sec:defaultpp_rq3}

DEFault++ separates root causes within predicted fault categories at Macro-F1 of 0.851 on encoders and 0.867 on decoders, as reported in Tables~\ref{tab:perf_level3_overall} and~\ref{tab:perf_level3_category}.

\begin{table}[!htbp]
\centering
\caption{DEFault++ performance on hierarchical root-cause diagnosis.}
\label{tab:perf_level3_overall}
\begin{tabular}{l c c}
\toprule
\textbf{Metric} & \textbf{Encoder} & \textbf{Decoder} \\
\midrule
Accuracy & 0.7953 $\pm$ 0.0189 & 0.8264 $\pm$ 0.0154 \\
Macro-F1 & 0.8512 $\pm$ 0.0160 & 0.8670 $\pm$ 0.0141 \\
Weighted-F1 & 0.8567 $\pm$ 0.0156 & 0.8745 $\pm$ 0.0134 \\
Precision & 0.9407 $\pm$ 0.0078 & 0.9475 $\pm$ 0.0154 \\
Recall & 0.7932 $\pm$ 0.0182 & 0.8225 $\pm$ 0.0136 \\
\bottomrule
\end{tabular}
\end{table}

\begin{table}[!htbp]
\centering
\caption{Per-category root-cause performance for DEFault++.}
\label{tab:perf_level3_category}
\small
\begin{tabular}{l c c c c c c c c}
\toprule
& \multicolumn{4}{c}{\textbf{Encoder}} & \multicolumn{4}{c}{\textbf{Decoder}} \\
\cmidrule(lr){2-5} \cmidrule(lr){6-9}
\textbf{Category} & \textbf{F1} & \textbf{Acc.} & \textbf{Prec.} & \textbf{Rec.} & \textbf{F1} & \textbf{Acc.} & \textbf{Prec.} & \textbf{Rec.} \\
\midrule
Embedding & 0.879 & 0.805 & 0.972 & 0.806 & 0.877 & 0.890 & 0.972 & 0.816 \\
FFN & 0.804 & 0.808 & 0.847 & 0.778 & 0.883 & 0.777 & 1.000 & 0.803 \\
Kernel & 0.814 & 0.740 & 1.000 & 0.706 & 0.747 & 0.691 & 0.988 & 0.676 \\
KV Cache & --- & --- & --- & --- & 0.881 & 0.822 & 1.000 & 0.817 \\
LayerNorm & 0.821 & 0.748 & 0.917 & 0.749 & 0.890 & 0.749 & 0.969 & 0.839 \\
Masking & 0.925 & 0.871 & 0.986 & 0.874 & 0.845 & 0.755 & 0.902 & 0.801 \\
Output & 0.882 & 0.825 & 0.958 & 0.822 & 0.972 & 0.969 & 1.000 & 0.947 \\
Positional & 0.930 & 0.884 & 0.992 & 0.886 & 0.975 & 0.911 & 1.000 & 0.953 \\
QKV & 0.765 & 0.693 & 0.868 & 0.694 & 0.646 & 0.629 & 0.727 & 0.592 \\
Residual & 0.908 & 0.880 & 0.927 & 0.895 & 0.958 & 0.955 & 0.961 & 0.955 \\
Score & 0.835 & 0.779 & 0.879 & 0.799 & 0.908 & 0.897 & 0.926 & 0.897 \\
Variant & 0.799 & 0.717 & 1.000 & 0.717 & 0.824 & 0.871 & 0.925 & 0.773 \\
\midrule
\textbf{Macro avg.} & 0.851 & 0.795 & 0.941 & 0.793 & 0.867 & 0.826 & 0.948 & 0.823 \\
\bottomrule
\end{tabular}
\end{table}

We identify Positional (0.930 encoder, 0.975 decoder), Masking (0.925 encoder, 0.845 decoder), and Residual (0.908 encoder, 0.958 decoder) as the strongest categories. Positional root causes produce distinct patterns in attention entropy and positional score features. Residual root causes are distinguished because skip-connection integrity and gradient-norm features produce non-overlapping distributions across root causes within that category. QKV is the weakest category, with F1 of 0.765 on encoders and 0.646 on decoders. We find that QKV root causes activate overlapping feature groups because multiple root causes within the category affect projection weights, alignment measurements, and update flow measurements simultaneously. The shared activation patterns limit the margin the prototype-matching loss can enforce between root causes.

The per-category scores in Table~\ref{tab:perf_level3_category} measure how well Level~3 separates root causes once the fault category is fixed. They complement the overall hierarchical root-cause metric in Table~\ref{tab:perf_level3_overall}. Level~2 errors still constrain Level~3 performance because DEFault++ predicts root causes only within the predicted fault category. We observe that precision exceeds recall across all categories, with macro averages of 0.941 on encoders and 0.948 on decoders against recall of 0.793 and 0.823, respectively. We interpret this gap as evidence that DEFault++ is conservative at the root-cause level. When DEFault++ commits to a root cause, the prediction is almost always correct. It does, however, miss some true root-cause instances, most often in categories where root causes share overlapping feature group patterns.

We examine prototype-level training dynamics in Figures~\ref{fig:defaultpp_prototype_faithfulness}--\ref{fig:defaultpp_prototype_utilization}. In Figure~\ref{fig:defaultpp_prototype_faithfulness}, we find that both the cross-entropy objective and prototype agreement increase steadily over training, as does root-cause proto F1. We observe that the two optimization objectives align rather than compete. Gradient directions from the two losses remain consistent throughout training. In Figure~\ref{fig:defaultpp_prototype_margin}, we observe that prototype margin grows throughout training. The widening gap between the nearest correct and nearest incorrect prototype measures improved discrimination at the root-cause level. In Figure~\ref{fig:defaultpp_prototype_utilization}, we find that usage entropy remains high and the active prototype count stays near its maximum throughout training, confirming that no prototype collapse occurs. Without broad prototype utilization, the model would misdiagnose rare root causes by assigning them to majority prototypes.

\begin{figure}[!htbp]
\centering
\includegraphics[width=0.9\linewidth]{fig_prototype_faithfulness.png}
\caption{Prototype-label agreement during training.}
\label{fig:defaultpp_prototype_faithfulness}
\end{figure}

\begin{figure}[!htbp]
\centering
\includegraphics[width=0.9\linewidth]{fig_prototype_margin.png}
\caption{Prototype margin during training.}
\label{fig:defaultpp_prototype_margin}
\end{figure}

\begin{figure}[!htbp]
\centering
\includegraphics[width=0.9\linewidth]{fig_prototype_utilization.png}
\caption{Prototype utilization during training.}
\label{fig:defaultpp_prototype_utilization}
\end{figure}

\FloatBarrier

We validate the distance-based explanation through group ablation. For each sample, we identify the two feature groups contributing most to the prototype distance margin between the nearest correct and nearest incorrect prototype (Equation~\ref{eq:proto_distance}). We then zero out these groups and compare the resulting margin drop with the drop obtained by zeroing two randomly chosen groups. In Figure~\ref{fig:rq3_group_ablation}, we find that removing the highest-distance groups produces a larger margin drop than removing random groups across all evaluated categories on both architectures. We conclude that the highest-distance feature groups causally contribute to the margin between correct and incorrect prototypes. The ablation measures explanation faithfulness through margin sensitivity. Ground-truth fault locations are not available for validation.

\begin{figure}[!htbp]
\centering
\includegraphics[width=0.9\linewidth]{fig_RQ3_group_ablation_faithfulness.png}
\caption{Group ablation analysis for distance-based explanation.}
\label{fig:rq3_group_ablation}
\end{figure}

% Old importance table and overlay/importance figures removed (superseded by group ablation analysis).

\FloatBarrier

\begin{rqsummary}
\textbf{Summary of RQ$\mathbf{_3}$:} Root-cause Macro-F1 reaches 0.851 on encoders and 0.867 on decoders, with precision above 0.94. Level~2 errors constrain downstream root-cause accuracy. The group ablation confirms that the highest-distance feature groups are behaviorally relevant to the diagnosis.
\end{rqsummary}

%==============================================================================
\section{Case Study}
\label{sec:defaultpp_casestudy}

This case study applies DEFault++ to the QKV fusion fault from HuggingFace Diffusers Issue~\#11903. The run is illustrative, not a standalone benchmark.

\subsection{Description of Fault}

The faulty program in Listing~\ref{lst:qkv_fusion_bug} defines separate \texttt{q\_proj}, \texttt{k\_proj}, and \texttt{v\_proj} layers and then fuses them into \texttt{to\_qkv}. The code does not delete the original projection layers after fusion. The \texttt{forward} method uses only \texttt{to\_qkv}. LoRA is configured to attach to \texttt{q\_proj} and \texttt{v\_proj}. LoRA therefore adapts parameters that no longer participate in the computation graph.

The model trains without raising an exception and produces no NaN or Inf values. The loss curve decreases plausibly. Training metrics alone do not identify the faulty component. The observable symptom, unchanged output after fine-tuning, does not identify which projection layer is at fault.

\subsection{Outcome of DEFault++}

We fine-tuned the model using the instrumentation described in Section~\ref{sec:defaultpp_features}. We collected dynamic features, including projection update ratios, attention entropy, and score statistics. We then ran DEFault++ on the resulting trace without supplying ground-truth fault information. Across training, the stale \texttt{q\_proj} and \texttt{v\_proj} layers produced near-zero update ratios. The fused projection layer continued to participate in the forward pass. This mismatch concentrated deviations in the QKV feature groups.

Figure~\ref{fig:defaultpp_case_qkv} reports the hierarchical predictions and the normalized per-group salience scores ($w_g$ from Equation~\ref{eq:explanation}) for a representative inference from this run.

\begin{figure}[!htbp]
\centering
\begin{tikzpicture}[
    arrow/.style={-{Stealth[length=5pt]}, thick, black!50},
    barfill/.style={fill=rqblue, draw=rqblue!80},
    barfilllight/.style={fill=rqblue!25, draw=rqblue!40},
]

% === Layout constants ===
\def\labelwidth{3.6}   % width of level-label column
\def\contentstart{3.9}  % x where prediction text begins
\def\boxleft{0.0}
\def\boxright{13.0}
\def\boxmid{6.5}

% === Panel (a): Hierarchical diagnosis ===
\node[font=\small\bfseries, anchor=west] at (\boxleft, 0.3) {(a) Hierarchical diagnostic output};

% --- Level 1 ---
\draw[black!50, rounded corners=3pt] (\boxleft, -0.2) rectangle (\boxright, -1.4);
\fill[rqblue, rounded corners=2pt] (\boxleft+0.15, -0.5) rectangle (\boxleft+\labelwidth, -1.1);
\node[text=white, font=\footnotesize\bfseries] at (\boxleft+\labelwidth/2+0.075, -0.8) {Level 1: Detection};
\node[anchor=west, font=\small] at (\contentstart, -0.55) {\textbf{Faulty} \quad (confidence: 0.98)};
\node[anchor=west, font=\scriptsize, text=black!60] at (\contentstart, -1.05) {Normal training loss, abnormal projection-update behavior};

% Arrow
\draw[arrow] (\boxmid, -1.4) -- (\boxmid, -1.7);

% --- Level 2 ---
\draw[black!50, rounded corners=3pt] (\boxleft, -1.7) rectangle (\boxright, -2.9);
\fill[rqblue, rounded corners=2pt] (\boxleft+0.15, -2.0) rectangle (\boxleft+\labelwidth, -2.6);
\node[text=white, font=\footnotesize\bfseries] at (\boxleft+\labelwidth/2+0.075, -2.3) {Level 2: Categorization};
\node[anchor=west, font=\small] at (\contentstart, -2.05) {\textbf{QKV} fault category \quad (confidence: 0.91)};
\node[anchor=west, font=\scriptsize, text=black!60] at (\contentstart, -2.55) {Deviations in QKV alignment, projection update, attention entropy};

% Arrow
\draw[arrow] (\boxmid, -2.9) -- (\boxmid, -3.2);

% --- Level 3 ---
\draw[black!50, rounded corners=3pt] (\boxleft, -3.2) rectangle (\boxright, -4.4);
\fill[rqblue, rounded corners=2pt] (\boxleft+0.15, -3.5) rectangle (\boxleft+\labelwidth, -4.1);
\node[text=white, font=\footnotesize\bfseries] at (\boxleft+\labelwidth/2+0.075, -3.8) {Level 3: Root cause};
\node[anchor=west, font=\small] at (\contentstart, -3.55) {\textbf{Stale-parameter fault} \quad (confidence: 0.88)};
\node[anchor=west, font=\scriptsize, text=black!60] at (\contentstart, -4.05) {Near-zero update ratio for \texttt{q\_proj}/\texttt{v\_proj}, fused path active};

% === Panel (b): Salience bar chart ===
\node[font=\small\bfseries, anchor=west] at (\boxleft, -5.2) {(b) Feature-group salience ($w_g$)};

\def\barlabelwidth{2.8}  % width reserved for bar labels
\def\barstart{3.0}
\def\barscale{7.0}
\def\barheight{0.32}
\def\bartextoffset{0.15}

% QKV alignment: 43%
\node[anchor=east, font=\scriptsize] at (\barstart - 0.15, -5.9) {QKV alignment};
\fill[barfill] (\barstart, -5.9-\barheight/2) rectangle (\barstart+0.43*\barscale, -5.9+\barheight/2);
\node[anchor=west, font=\scriptsize\bfseries] at (\barstart+0.43*\barscale+\bartextoffset, -5.9) {43\%};

% Projection update: 29%
\node[anchor=east, font=\scriptsize] at (\barstart - 0.15, -6.4) {Projection update};
\fill[barfill] (\barstart, -6.4-\barheight/2) rectangle (\barstart+0.29*\barscale, -6.4+\barheight/2);
\node[anchor=west, font=\scriptsize\bfseries] at (\barstart+0.29*\barscale+\bartextoffset, -6.4) {29\%};

% Attention statistics: 13%
\node[anchor=east, font=\scriptsize] at (\barstart - 0.15, -6.9) {Attention statistics};
\fill[barfilllight] (\barstart, -6.9-\barheight/2) rectangle (\barstart+0.13*\barscale, -6.9+\barheight/2);
\node[anchor=west, font=\scriptsize] at (\barstart+0.13*\barscale+\bartextoffset, -6.9) {13\%};

% Score statistics: 7%
\node[anchor=east, font=\scriptsize] at (\barstart - 0.15, -7.4) {Score statistics};
\fill[barfilllight] (\barstart, -7.4-\barheight/2) rectangle (\barstart+0.07*\barscale, -7.4+\barheight/2);
\node[anchor=west, font=\scriptsize] at (\barstart+0.07*\barscale+\bartextoffset, -7.4) {7\%};

% Training dynamics: 5%
\node[anchor=east, font=\scriptsize] at (\barstart - 0.15, -7.9) {Training dynamics};
\fill[barfilllight] (\barstart, -7.9-\barheight/2) rectangle (\barstart+0.05*\barscale, -7.9+\barheight/2);
\node[anchor=west, font=\scriptsize] at (\barstart+0.05*\barscale+\bartextoffset, -7.9) {5\%};

% Other groups: 3%
\node[anchor=east, font=\scriptsize] at (\barstart - 0.15, -8.4) {Other groups};
\fill[barfilllight] (\barstart, -8.4-\barheight/2) rectangle (\barstart+0.03*\barscale, -8.4+\barheight/2);
\node[anchor=west, font=\scriptsize] at (\barstart+0.03*\barscale+\bartextoffset, -8.4) {3\%};

% Annotation brace for projection-related groups
\draw[decorate, decoration={brace, amplitude=5pt, mirror}, thick, rqblue]
    (\barstart+0.43*\barscale+0.8, -5.9+\barheight/2+0.05)
    -- (\barstart+0.43*\barscale+0.8, -6.4-\barheight/2-0.05)
    node[midway, right=6pt, font=\scriptsize, text=rqblue, align=left] {72\% in projection-\\related groups};

\end{tikzpicture}
\caption{DEFault++ diagnostic output for the stale QKV fusion fault (Listing~\ref{lst:qkv_fusion_bug}). Panel~(a) shows the hierarchical predictions with confidence scores at each level. Panel~(b) shows the normalized feature-group salience ($w_g$) from Equation~\ref{eq:explanation}. The two projection-related groups account for 72\% of the total salience.}
\label{fig:defaultpp_case_qkv}
\end{figure}

\paragraph{Level 1: Fault detection.}
DEFault++ applies binary detection to the trace. In the representative inference summarized in Figure~\ref{fig:defaultpp_case_qkv}, the detector assigns a high fault score to the run. Projection-update features drive the detection, not global training failure. Loss decreases normally. The internal projection-update measurements deviate from the expected behavior of a correctly functioning attention block. Training monitors that observe only loss curves and gradient norms do not detect this pattern. The optimization process produces no anomalous aggregate metric.

\paragraph{Level 2: Fault categorization.}
DEFault++ assigns the trace to the \textit{QKV} fault category. The category-localizing deviations concentrate in three feature groups. First, the projection update ratios for \texttt{q\_proj} and \texttt{v\_proj} remain near zero, indicating no effective update. Second, the pairwise Q--K, Q--V, and K--V alignment features deviate from the pattern of actively fine-tuned projections. Third, attention entropy and score statistics shift toward patterns associated with disturbed projection coherence, not masking, positional, or residual-stream faults. DEFault++ outputs \textit{QKV} as the fault category and forwards that label to Level~3.

\paragraph{Level 3: Root cause diagnosis.}
DEFault++ predicts a \textit{weight-update / stale-parameter fault}. DEFault++ reaches that prediction by comparing the trace with per-root-cause prototypes in the embedding space. In the representative explanation, the largest normalized distance shares fall on QKV alignment and projection update dynamics. Task metrics and global training features receive smaller shares. That distribution corresponds to the stale-parameter scenario. The model trains, but the adapted parameters are disconnected from the computation that produces the output.

The code-level evidence for that explanation appears in Listing~\ref{lst:qkv_fusion_bug}. The \texttt{forward} method consumes \texttt{to\_qkv}. The \texttt{q\_proj} and \texttt{v\_proj} attributes remain addressable to downstream adapters. The distance-based salience concentrates on the QKV and projection-update groups, corresponding to that code path. We did not perform independent developer validation of this explanation. To confirm the explanation, verify that \texttt{q\_proj}, \texttt{k\_proj}, and \texttt{v\_proj} remain reachable after fusion and that LoRA targets the fused layer, not the stale module attributes.

AutoTrainer~\cite{autotrainer} monitors training health metrics such as loss trends and gradient behavior. The global optimization in this run appears healthy by those metrics. AutoTrainer does not flag a fault. DeepDiagnosis~\cite{deepdiagnosis} uses aggregate training metrics and does not measure projection-level update behavior. DeepFD~\cite{deepfd} learns from runtime features. Its label space covers general DNN fault categories and does not include transformer-specific projection mechanisms such as stale post-fusion layers. DeepLocalize~\cite{deeplocalize} detects numerical errors such as NaN or infinity values in loss or activations. This run produces neither. DEFault++ uses transformer-specific features that directly measure projection alignment and update behavior. It decomposes the diagnosis into per-component distance shares that concentrate on the projection-related feature groups.


%==============================================================================
\section{Threats to Validity}
\label{sec:defaultpp_threats}

Threats to \emph{internal validity} relate to experimental errors. The main threat is information leakage among configurations sharing the same model, dataset, and seed. We mitigate this through nested grouped cross-validation, fitting preprocessing only on the inner-training fold, and fixing the outer fold assignment across all compared methods. A second threat is the joint training objective, where the shared feature-encoding pathway means that changes in one loss can affect the representation available to the others. We partially examine this interaction through the ablation in RQ$\mathbf{_2}$.

Threats to \emph{construct validity} relate to whether the benchmark captures the intended diagnosis problem. The benchmark is produced by mutation operators rather than by naturally occurring development histories~\cite{attention_taxonomy,frankenformer}. Five-seed averaging may also obscure transient behaviors, and the benchmark provides limited coverage of compound or co-occurring faults. Class imbalance is another construct-validity threat. We addressed this through Macro-F1, inverse-frequency weighting, and oversampling. The explanation mechanism is a further construct-validity concern. The distance-based salience quantities are model-internal and do not represent causal effects. We validate them through group ablation (RQ$\mathbf{_3}$) and a code-level check in the case study, but they remain model-internal quantities rather than externally validated causal accounts.

Threats to \emph{external validity} relate to generalizability. The dataset covers seven transformer models and nine tasks~\cite{frankenformer}, but results may change for much larger models or architectural variants not represented in the benchmark (e.g., mixture-of-experts, multimodal cross-attention). The FPG is derived from standard transformer blocks, so materially different architectures may require a different propagation graph. Additionally, existing baselines do not expose a label space compatible with our fault categories, so direct comparison is limited to Level~1 detection. Five grouped outer folds provide a limited number of independent test partitions. We report per-fold standard deviations and use the same folds across compared methods to make variation visible.


%==============================================================================
\section{Related Work}
\label{sec:defaultpp_related}

\textbf{Debugging DNN Programs.}
Rule-based debuggers such as AutoTrainer~\cite{autotrainer}, DeepDiagnosis~\cite{deepdiagnosis}, and UMLAUT~\cite{umlaut} reason over generic training symptoms. Learning-based approaches such as DeepFD~\cite{deepfd} and our prior work DEFault~\cite{default} move from fixed rules to classifiers over runtime traces. These approaches are most closely aligned with DEFault++, but they still predict DNN-level fault types rather than transformer component categories and within-category root causes. For this chapter, that distinction matters more than the choice of classifier. A label space built for loss functions, activations, or optimizers does not directly express QKV, masking, positional, residual, or KV-cache faults. DeepLocalize~\cite{deeplocalize} traces numerical failures back to layers from runtime values, while static analyzers such as NeuraLint~\cite{NeuraLint}, NerdBug~\cite{nerdbug}, DEBAR~\cite{debar}, and BugLab~\cite{buglab} inspect structure without training execution. FL4Deep~\cite{fl4deep} combines static code features with dynamic features through a knowledge graph. DEFault++ instead targets runs that remain structurally valid and numerically stable but leave component-specific behavioral traces during training.

\textbf{Transformer-Specific Analysis.}
Some prior work is transformer-specific but addresses a different fault model. Adversarial evaluation, behavioral testing, and attention analysis study robustness and behavior under perturbation~\cite{jones2020adversarial,ribeiro2020checklist,clark2019does,voita2019analyzing}. Systems-oriented methods such as ATTNChecker~\cite{liang2025attnchecker}, FT-Transformer~\cite{fttransformer}, AtPatch~\cite{atpatch2026}, and Bug Attention Probe~\cite{bap2025} focus on transient numerical faults, tolerance mechanisms, or attention-map anomalies. These studies show that transformers require dedicated analysis, but they do not address supervised diagnosis of persistent implementation faults in a component/root-cause label space. LLM-based debugging assistants such as SoapFL~\cite{soapfl2025}, LLM4FL~\cite{llm4fl2024}, ChatDBG~\cite{chatdbg2025}, and BugReAct~\cite{bugreact2026} assist source-level debugging in general software rather than learning a transformer fault classifier from instrumented traces.

\textbf{Mutation Testing.}
Supervised fault diagnosis depends on labeled faulty runs, and mutation testing has been the main way to construct such data. DeepMutation~\cite{deepmutation}, DeepMutation++~\cite{deepmutation++}, and DeepCrime~\cite{deepcrime} expanded mutation testing for DNNs, but their operators target generic DNN components rather than attention-specific mechanisms. FrankenFormer~\cite{frankenformer} injects transformer-specific faults and records internal measurements relevant to attention behavior. DEFault++ uses that benchmark as the training and evaluation basis, which is why its claims should be read in the context of mutation-based supervision rather than unconstrained field data.

\textbf{Explainability.}
Most explanation methods used in software engineering and prior fault diagnosis are post-hoc. Representative examples include SHAP, DiCE, and Anchors~\cite{lundberg2017unified,mothilal2020dice,ribeiro2018anchors}, which return feature attributions, counterfactuals, or rules after the classifier has already made its decision. DEFault (Chapter~\ref{chapter_default}) used this style of explanation in a cascaded way for diagnosis. DEFault++ uses a different strategy. The Level~3 diagnosis itself is defined through prototype comparison in a grouped embedding space, and the reported explanation is a decomposition of that same prototype distance across feature groups. This design is related to prototype networks~\cite{snell2017prototypical} and concept-bottleneck ideas~\cite{koh2020concept}, but here the groups correspond to transformer subsystems rather than to human semantic concepts.


%==============================================================================
\section{Contribution}
\label{sec:defaultpp_summary}

This chapter introduced DEFault++, a transformer-specific extension of DEFault. DEFault showed that hierarchical diagnosis can support fault detection and categorization for FFNNs, CNNs, and RNNs using DNN-level static and dynamic features. DEFault++ carries that hierarchical idea into the transformer setting, but changes the objects of diagnosis. It uses transformer-specific runtime measurements, predicts component-level fault categories and within-category root causes, and organizes the diagnosis around a Fault Propagation Graph so that the returned evidence remains tied to transformer subsystems. On the FrankenFormer benchmark, DEFault++ achieves AUROC above 0.96 for detection, Macro-F1 above 0.85 for categorization, and hierarchical root-cause Macro-F1 above 0.85 on both encoder and decoder architectures. These results show that transformer traces are informative enough to support automated fault detection, categorization, and hierarchical root-cause diagnosis. Future work can focus on generalizing the proposed hierarchical technique to other architectures with well-defined component structures (e.g., mixture-of-experts, multimodal cross-attention models) by adapting the Fault Propagation Graph (FPG) and corresponding feature groups directly to the target architecture. 